\documentclass[12pt,a4paper]{article}

\usepackage[T2A]{fontenc}
\usepackage[utf8]{inputenc}
\usepackage[russian]{babel}

\usepackage{amsmath,amssymb,amsthm}
\usepackage{geometry}
\usepackage{hyperref}
\usepackage{microtype}
\usepackage{setspace}
\usepackage{caption}
\usepackage{graphicx}

\geometry{
  left=25mm,
  right=25mm,
  top=25mm,
  bottom=25mm
}

\hypersetup{
  colorlinks=true,
  linkcolor=black,
  urlcolor=blue
}

\begin{document}
\selectlanguage{russian}

% ================= ТИТУЛЬНЫЙ ЛИСТ =================

\begin{titlepage}
  \begin{center}
    {Санкт-Петербургский политехнический университет Петра Великого\\
    Институт прикладной математики и механики\\[0.5em]
    \textbf{Высшая школа прикладной математики и вычислительной физики}}

    \vfill

    {\LARGE\bfseries Отчёт по курсовой работе\\ Часть №2\par}
    \vspace{1em}
    {\Large по дисциплине\\[0.3em]
    {\bfseries «Проект по математической статистике»}}

    \vfill

    \begin{flushright}
        Выполнил\\
        студент гр. 5030102/20202:\\
        Зайдиев Артур
    \end{flushright}

    \begin{flushright}
        Проверил\\
        Преподаватель: Заяц Олег Иванович
    \end{flushright}

    \vfill
    Санкт-Петербург\\
    2026
  \end{center}
\end{titlepage}

\setcounter{page}{2}
\newpage

% ================= ОСНОВНАЯ ЧАСТЬ =================

\section*{Вычисление ОИСКО ядерной оценки: показательное распределение и колокольное ядро}

Задана плотность показательного распределения:

\[
f_0(x)=
\begin{cases}
\lambda e^{-\lambda x}, & x \ge 0,\\
0, & x<0
\end{cases}
\]

где $\lambda = 1$.

Для оценки плотности используется ядерная оценка с колокольным (гауссовским) ядром:

\[
K(u)=\frac{1}{\sqrt{2\pi}}e^{-u^2/2}, \quad \hat{f}_n(x) = \frac{1}{n\xi}\sum_{i=1}^n K\left( \frac{x - X_i}{\xi} \right).
\]

Относительное интегральное среднеквадратическое отклонение (ОИСКО) вычисляется по формуле:

\[
\bar{\delta}_n(\xi)=
1+\frac{I_1}{n\xi I_4}
+\left(1-\frac{1}{n}\right)\frac{I_2(\xi)}{I_4}
-2\frac{I_3(\xi)}{I_4}.
\]

\subsection*{Вычисление интегралов}

Были получены следующие аналитические выражения:

\[
I_1 = \sqrt{\pi}, \quad I_4 = \pi\lambda = \pi,
\]

\[
I_2(\xi) = \pi\lambda \, e^{\lambda^2 \xi^2} \operatorname{erfc}(\lambda \xi),
\]

\[
I_3(\xi) = \pi\lambda \, e^{\lambda^2 \xi^2 / 2} \operatorname{erfc}\left( \frac{\lambda \xi}{\sqrt{2}} \right).
\]

При $\lambda=1$ окончательная формула ОИСКО принимает вид:

\[
\bar{\delta}_n(\xi)=
1 + \frac{\sqrt{\pi}}{n\xi \pi}
+ \left(1-\frac{1}{n}\right) e^{\xi^2} \operatorname{erfc}(\xi)
- 2 e^{\xi^2 / 2} \operatorname{erfc}\left( \frac{\xi}{\sqrt{2}} \right).
\]

Оптимальное значение параметра сглаживания $\xi^*$ находится путём численной минимизации $\bar{\delta}_n(\xi)$.

\newpage

\section*{Построение графиков}

Для анализа качества ядерной оценки и сравнения её эффективности с гистограммной оценкой были построены следующие графики.

\subsection*{График 1. Зависимость ОИСКО $\bar{\delta}_n(\xi)$ от параметра $\xi$}

\begin{figure}[ht!]
    \centering
    \includegraphics[width=0.85\textwidth]{figures/f1_kernel.png}
    \caption{ОИСКО ядерной оценки $\bar{\delta}_n(\xi)$ для $n=10,50,250,1000$}
    \label{fig:kernel_delta}
\end{figure}

\subsection*{График 2. Зависимость оптимального параметра $\xi_{opt}$ от объёма выборки $n$}

\begin{figure}[ht!]
    \centering
    \includegraphics[width=0.85\textwidth]{figures/f2_xi_opt.png}
    \caption{Оптимальный параметр сглаживания $\xi_{opt}$ ядерной оценки}
    \label{fig:xi_opt}
\end{figure}

\newpage

\subsection*{График 3. Зависимость нижней границы погрешности $\bar{\delta}_{n,\min}$ от $n$}

\begin{figure}[ht!]
    \centering
    \includegraphics[width=0.85\textwidth]{figures/f3_delta_min.png}
    \caption{Минимальное значение ОИСКО $\bar{\delta}_{n,\min}$ ядерной оценки}
    \label{fig:delta_min}
\end{figure}

\subsection*{График 4. Сравнение ядерной оценки и гистограммы}

\begin{figure}[ht!]
    \centering
    \includegraphics[width=0.85\textwidth]{figures/f4_comparison.png}
    \caption{Сравнение минимального ОИСКО ядерной оценки (колокольное ядро) и гистограммной оценки}
    \label{fig:comparison}
\end{figure}

\end{document}